% Блок-схема алгоритма пересчёта результатов отбора.
% Оформление по ГОСТ 19.701-90:
%   прямоугольник — операция/процесс;
%   ромб — условие;
%   овал-«пилюля» — начало/конец;
%   параллелограмм — ввод/вывод данных.
%
% Документ компилируется самостоятельно через XeLaTeX:
%   xelatex selection-algorithm.tex
% Результат — selection-algorithm.pdf одностраничный, размер кадрируется по содержимому.

\documentclass[border=4mm]{standalone}

\usepackage{fontspec}
\setmainfont{Times New Roman}

\usepackage{tikz}
\usetikzlibrary{shapes.geometric, arrows.meta, positioning, calc}

\begin{document}

\begin{tikzpicture}[
  font=\fontsize{10pt}{12pt}\selectfont,
  node distance=6mm and 12mm,
  every node/.style={align=center},
  % Стили блоков ГОСТ 19.701-90
  startstop/.style={
    rectangle, rounded corners=4mm,
    minimum width=46mm, minimum height=9mm,
    draw=black, thick, fill=white,
  },
  process/.style={
    rectangle,
    minimum width=70mm, minimum height=10mm,
    draw=black, thick, fill=white,
    text width=66mm,
  },
  io/.style={
    trapezium, trapezium left angle=70, trapezium right angle=110,
    minimum width=70mm, minimum height=9mm,
    draw=black, thick, fill=white,
    text width=64mm,
  },
  decision/.style={
    diamond, aspect=2.8,
    minimum width=74mm, minimum height=15mm,
    draw=black, thick, fill=white,
    inner sep=1pt,
    text width=60mm,
  },
  arrow/.style={-{Latex[length=2.5mm, width=2mm]}, thick},
  label/.style={font=\fontsize{9pt}{11pt}\selectfont, inner sep=3pt, outer sep=3pt, fill=white},
]

% ===== Основной поток =====
\node[startstop]                       (start)    {Начало};
\node[io,       below=of start]        (loaddata) {Загрузить снимок данных из БД};
\node[process,  below=of loaddata]     (filter)   {Исключить заявки с незаполненными\\ критериями группы категории};
\node[process,  below=of filter]       (init)     {Поместить всех абитуриентов\\ с непустыми списками заявок в очередь};

\node[decision, below=of init]         (qempty)   {Очередь пуста?};

\node[process,  below=of qempty]       (take)     {Извлечь абитуриента из очереди};
\node[decision, below=of take]         (haspref)  {Есть категория\\ в списке предпочтений?};
\node[process,  below=of haspref]      (addpool)  {Добавить в пул кандидатов\\ следующей категории};

\node[decision, below=of addpool]      (overquota) {Размер пула\\ больше квоты?};
\node[process,  below=of overquota]    (evict1)   {Извлечь самого слабого\\ кандидата и вернуть в очередь};

\node[decision, below=of evict1]       (overplan)  {Сумма зачисленных\\ больше плана?};
\node[process,  below=of overplan]     (evict2)   {Извлечь самого слабого из категории с наименьшим приоритетом};

% Финальные блоки — со сдвигом вниз, чтобы освободить место под loopback
\node[process,  below=of evict2, yshift=-10mm]   (clear)   {Очистить таблицу результатов};
\node[process,  below=of clear]                  (insert)  {Массово вставить содержимое пулов};
\node[startstop, below=of insert]                (stop)    {Конец};

% Блок отказа — справа от ромба «есть предпочтения»
\node[process, right=30mm of haspref] (reject) {Абитуриент не принят ни по одной из категорий};

% ===== Связи =====
\draw[arrow] (start) -- (loaddata);
\draw[arrow] (loaddata) -- (filter);
\draw[arrow] (filter) -- (init);
\draw[arrow] (init) -- (qempty);

% Основной нисходящий поток с подписями условий
\draw[arrow] (qempty) -- node[label, right] {Нет} (take);
\draw[arrow] (take) -- (haspref);
\draw[arrow] (haspref) -- node[label, right] {Да} (addpool);
\draw[arrow] (addpool) -- (overquota);

\draw[arrow] (overquota) -- node[label, right] {Да} (evict1);
\draw[arrow] (evict1) -- (overplan);
\draw[arrow] (overplan) -- node[label, right] {Да} (evict2);

% Ветка «Нет» от overquota — обход справа, минуя evict1, сразу к проверке плана
\draw[arrow] (overquota.east) -- node[label, pos=0, anchor=west, above, xshift=1pt] {Нет} ++(15mm,0) |- (overplan.east);

% Ветка «Нет» от haspref — в блок отказа
\draw[arrow] (haspref.east) -- node[label, pos=0, anchor=west, above, xshift=1pt] {Нет} (reject.west);

% Возврат в очередь: evict2, overplan-Нет и reject сходятся к левому спайну,
% который поднимается до qempty.west
\coordinate (loopback-start)   at ($(evict2.south)+(0,-4mm)$);
\coordinate (loopback-corner1) at ($(evict2.south west)+(-25mm,-4mm)$);
\coordinate (loopback-corner2) at ($(qempty.west)+(-25mm,0)$);
% Точка на левом спайне, на уровне overplan — для ветки «Нет» от overplan
\coordinate (loopback-overplan) at (loopback-corner1 |- overplan.west);

\draw[arrow] (evict2.south) -- (loopback-start) -- (loopback-corner1) -- (loopback-corner2) -- (qempty.west);
\draw      (reject.south) |- (loopback-start);
% «Нет» от overplan: план не превышен, вытеснение не нужно — возврат в очередь
\draw      (overplan.west) -- node[label, midway, above] {Нет} (loopback-overplan);

% Ветка «Да» от qempty — выход из цикла к блоку очистки.
% Идёт справа, чтобы не пересекаться с loopback слева.
\draw[arrow] (qempty.east) -- node[label, pos=0, anchor=west, above, xshift=1pt] {Да} ++(20mm,0) |- (clear.east);

\draw[arrow] (clear) -- (insert);
\draw[arrow] (insert) -- (stop);

\end{tikzpicture}

\end{document}
